% Generated draft: requires mathematical and source review.
\documentclass[11pt]{article}
\usepackage[margin=25mm]{geometry}
\usepackage{fontspec}
\setmainfont{Latin Modern Roman}
\usepackage{amsmath,amssymb,mathtools,unicode-math}
\setmathfont{Latin Modern Math}
\usepackage{enumitem,needspace,xurl,hyperref}
\setlength{\parindent}{0pt}
\setlength{\parskip}{5pt}
\setlength{\emergencystretch}{3em}
\newcommand{\N}{\mathbb{N}}
\newcommand{\Z}{\mathbb{Z}}
\newcommand{\Q}{\mathbb{Q}}
\newcommand{\R}{\mathbb{R}}
\newcommand{\C}{\mathbb{C}}
\newcommand{\F}{\mathbb{F}}
\newcommand{\A}{\mathbb{A}}
\newcommand{\PP}{\mathbb{P}}
\newcommand{\E}{\mathbb{E}}
\newcommand{\eps}{\varepsilon}
\newcommand{\dd}{\,\mathrm{d}}
\begin{document}
\section*{UnsolvedMath ID: 9400049}
\textbf{Source problem code:} catalogue entry
\subsection{Definitions and mathematical statement}
Fix an admissible local Gan--Gross--Prasad datum: a local field, the intended classical-group pair, its subgroup $H$, the auxiliary representation $\nu$, and a generic local $L$-parameter $\phi$. Let $\Pi_\phi^{\mathrm{rel}}$ denote the intended relevant packet, with its range of inner forms fixed. Let $\iota_\phi$ be its component-group labeling, with the necessary normalization data, and let $\chi_{\mathrm{GGP}}(\phi)$ be the distinguished Gan--Gross--Prasad character.

For an irreducible representation $\pi$ in this packet, the space
\[
\operatorname{Hom}_H(\pi\otimes\overline\nu,\mathbb C)
\]
consists of the appropriate smooth/continuous $H$-invariant linear forms after restriction to $H$. The conjecture asks for exactly one $\pi_0\in\Pi_\phi^{\mathrm{rel}}$ with this space nonzero, and identifies it by
\[
\iota_\phi(\pi_0)=\chi_{\mathrm{GGP}}(\phi).
\]
Equivalently, the subset of relevant packet members with a nonzero invariant-form space is the singleton consisting of the distinguished-character member. The source asks nonvanishing and uniqueness; no additional assertion about the dimension of a nonzero Hom space is inserted.

The actual datum, Bessel/Fourier--Jacobi case, local-field/characteristic range, packet/inner-form conventions, meaning of generic, and epsilon-factor character formula are not supplied. They are indispensable to an explicit case-by-case statement. The conjugate $\overline\nu$ is retained rather than replaced without conditions by a different representation. The supplied summary reports broad local resolutions by Waldspurger, Beuzart-Plessis, Gan, and others but gives no precise theorem hypotheses.
\subsection{Short English statement}
For each intended generic local GGP packet, prove that exactly one relevant representation has the prescribed invariant linear form, and that its label is the distinguished GGP character. The full local datum and normalization conventions are omitted.
\subsection{Sources}
\begin{itemize}
\item Wikipedia contributors, \emph{List of unsolved problems in mathematics}, Algebraic number theory bullet 49, pinned revision 1366636767, \url{https://en.wikipedia.org/w/index.php?title=List_of_unsolved_problems_in_mathematics&oldid=1366636767#Number_theory}; accessed 2026-07-29, with linked-page resolution recorded in the provenance. The question is attributed to Wee Teck Gan, Benedict Gross, and Dipendra Prasad.
\item The supplied summary refers to the Wikipedia page ``Gan--Gross--Prasad conjecture'' and names Waldspurger, Beuzart-Plessis, Gan, and others as proof attributions, without publication identifiers or complete hypotheses.
\item Local registry \texttt{data/unsolvedmath/problems.json}, permanent UnsolvedMath ID $9400049$, source code catalogue entry, \emph{Local Gan--Gross--Prasad conjecture}. Dataset attribution: Ulam AI and the UnsolvedMath Contributors.
\end{itemize}
\end{document}
